% Created 2025-02-07 Fri 13:23
% Intended LaTeX compiler: pdflatex
\documentclass[a4paper]{article}
\usepackage[utf8]{inputenc}
\usepackage[T1]{fontenc}
\usepackage{graphicx}
\usepackage{longtable}
\usepackage{wrapfig}
\usepackage{rotating}
\usepackage[normalem]{ulem}
\usepackage{amsmath}
\usepackage{amssymb}
\usepackage{capt-of}
\usepackage{hyperref}
\usepackage[spanish]{babel}
\usepackage[usenames,dvipsnames]{color} % Required for custom colors
\renewcommand{\ttdefault}{pcr} % MONOESPACIO CON NEGRIT
\usepackage{lastpage}
\usepackage{listings}
\usepackage{listingsutf8}
\renewcommand{\lstlistingname}{Listado}
\lstset{frame=single,inputencoding=utf8,basicstyle=\scriptsize\ttfamily,showstringspaces=false,numbers=none}
\definecolor{MyDarkGreen}{rgb}{0.0,0.4,0.0} % This is the color used for comments
\lstset{ breaklines=true, postbreak=\mbox{\textcolor{red}{$\hookrightarrow$}\space}, keywordstyle=\bfseries, keywordstyle=[1]\color{Blue}\bfseries,  keywordstyle=[2]\color{Purple}\bfseries,  keywordstyle=[3]\color{Blue}\underbar,   identifierstyle=,   commentstyle=\usefont{T1}{pcr}{m}{sl}\color{MyDarkGreen}\small,   stringstyle=\color{Purple},   showstringspaces=false,   tabsize=2,   morecomment=[l][\color{Blue}]{...} }
\lstset{literate=  {á}{{\'a}}1 {é}{{\'e}}1 {í}{{\'i}}1 {ó}{{\'o}}1 {ú}{{\'u}}1   {Á}{{\'A}}1 {É}{{\'E}}1 {Í}{{\'I}}1 {Ó}{{\'O}}1 {Ú}{{\'U}}1   {à}{{\`a}}1 {è}{{\`e}}1 {ì}{{\`i}}1 {ò}{{\`o}}1 {ù}{{\`u}}1   {À}{{\`A}}1 {È}{{\'E}}1 {Ì}{{\`I}}1 {Ò}{{\`O}}1 {Ù}{{\`U}}1   {ä}{{\"a}}1 {ë}{{\"e}}1 {ï}{{\"i}}1 {ö}{{\"o}}1 {ü}{{\"u}}1   {Ä}{{\"A}}1 {Ë}{{\"E}}1 {Ï}{{\"I}}1 {Ö}{{\"O}}1 {Ü}{{\"U}}1   {â}{{\^a}}1 {ê}{{\^e}}1 {î}{{\^i}}1 {ô}{{\^o}}1 {û}{{\^u}}1   {Â}{{\^A}}1 {Ê}{{\^E}}1 {Î}{{\^I}}1 {Ô}{{\^O}}1 {Û}{{\^U}}1   {œ}{{\oe}}1 {Œ}{{\OE}}1 {æ}{{\ae}}1 {Æ}{{\AE}}1 {ß}{{\ss}}1   {ű}{{\H{u}}}1 {Ű}{{\H{U}}}1 {ő}{{\H{o}}}1 {Ő}{{\H{O}}}1   {ç}{{\c c}}1 {Ç}{{\c C}}1 {ø}{{\o}}1 {å}{{\r a}}1 {Å}{{\r A}}1   {€}{{\euro}}1 {£}{{\pounds}}1 {«}{{\guillemotleft}}1   {»}{{\guillemotright}}1 {ñ}{{\~n}}1 {Ñ}{{\~N}}1 {¿}{{?`}}1 }
\usepackage{caption}
\usepackage{attachfile2}
\usepackage[margin=1.5cm,includeheadfoot,includehead,includefoot]{geometry}
\hypersetup{colorlinks,linkcolor=black}
\usepackage{fancyhdr}
\pagestyle{fancyplain}
\chead{}
\lhead{}
\rhead{}
\cfoot{}
\lfoot{\begin{footnotesize}alvaro.gonzalezsotillo@educa.madrid.org\end{footnotesize}}
\rfoot{\begin{footnotesize}\thepage / \pageref{LastPage}\end{footnotesize}}
\usepackage[skins]{tcolorbox}
\usepackage{multicol}
\usepackage{changepage} %ajdustwidth
\usepackage{fancybox}
\usepackage{attachfile2}
\lhead{Extraordinaria 2021 (es el \\lhead)}
\rhead{Administración y Gestión de Bases de Datos (es el \\rhead)}
\lhead{Trabajo de scripts de shell}
\rhead{ASGBD}
\usepackage{tabularx}
\usepackage{svg}
\usepackage{letltxmacro}
\LetLtxMacro{\originalincludegraphics}{\includegraphics}
\renewcommand{\includegraphics}[2][]{\IfFileExists{#2.pdf}{\originalincludegraphics[#1]{#2.pdf}}{\originalincludegraphics[#1]{#2}}}
\LetLtxMacro{\originalincludesvg}{\includesvg}
\renewcommand{\includesvg}[2][]{\IfFileExists{#2.pdf}{\originalincludegraphics[#1]{#2.pdf}}{\originalincludegraphics[#1]{#2.svg.pdf}}}
\usepackage{comment}
\excludecomment{NOTES}
\date{}
\title{}
\hypersetup{
 pdfauthor={Álvaro González Sotillo},
 pdftitle={},
 pdfkeywords={},
 pdfsubject={},
 pdfcreator={Emacs 29.4 (Org mode 9.6.15)}, 
 pdflang={Spanish}}
\begin{document}

\captionsetup{font=scriptsize}

\textattachfile{/home/alvaro/apuntes-clase-profesor/sistemas-gestores-bbdd-asir2/apuntes/3/asgbd-03-trabajo-scripts.org}{} \vspace{-1em}


\setlength{\parindent}{0em}
\setlength{\parskip}{1em}

\newtcolorbox{Aviso}[1][Aviso]{
  enhanced,
  colback=gray!5!white,
  colframe=gray!75!black,fonttitle=\bfseries,
  colbacktitle=gray!85!black,
  attach boxed title to top left={yshift=-2mm,xshift=2mm},
  title=#1
}

\newtcolorbox{cuadrito}[1][Ignorado]{
  %drop shadow southeast,
  enhanced jigsaw,
  colback=white,
}


\newcommand{\StudentData}{
  \begin{cuadrito}[1\textwidth]
    \vspace{0.3cm}
    \large{
      \textbf{Apellidos:} \hrulefill \\
      \textbf{Nombre:} \hrulefill \\
      \textbf{Fecha:} \hrulefill \hspace{1cm} \textbf{Usuario:} \hrulefill \\
    }
    \vspace{-0.2cm}
  \end{cuadrito}
}

\newcommand{\StudentDataDniFirma}{
  \begin{cuadrito}[1\textwidth]
    \vspace{0.3cm}
    \large{
      \textbf{Apellidos:} \hrulefill \\
      \textbf{Nombre:} \hrulefill  \\
      \textbf{Dni:} \hrulefill \hspace{8cm} \\
      \\
      \textbf{Firma:} \hrulefill \hspace{8cm} \\
    }
    \vspace{-0.2cm}
  \end{cuadrito}
}

La última versión de este documento se puede descargar de \url{https://alvarogonzalezsotillo.github.io/apuntes-clase/sistemas-gestores-bbdd-asir2/apuntes/3/asgbd-03-trabajo-scripts.pdf}




\section{\emph{Scripts} de inicio y parada de \textbf{Oracle} (1 punto)}
\label{sec:org0000000}

Crea dos \emph{scripts} para iniciar y parar \textbf{Oracle}
\begin{itemize}
\item \texttt{/home/alumno/scripts/oraclestart.sh}
\item \texttt{/home/alumno/scripts/oraclestop.sh}
\end{itemize}


\section{Arrancar automáticamente \textbf{Oracle} cuando se inicie el servidor (2 puntos)}
\label{sec:org0000003}


\begin{itemize}
\item \textbf{Oracle} debe levantarse cuando la máquina se inicie, y apagarse cuando la máquina se cierre.
\item Oracle se iniciará solo si se indica en el fichero \texttt{/etc/oratab}.

\item En el fichero \texttt{/home/alumno/logs/oracle.log} se dejará una traza de cuando se arrancó y se paró la máquina, y si fue necesario arrancar o parar \textbf{Oracle}. Por ejemplo:

\begin{lstlisting}[language=bash,label= ,caption={Ejemplo de \texttt{/home/alumno/logs/oracle.log} cuando \textbf{Oracle} se arranca},captionpos=b,numbers=none]
      2017-02-10-12:40:00 - Solicitud de arrancar Oracle
      2017-02-10-12:40:01 - Oracle arrancando porque /etc/oratab indica Y
      2017-02-10-12:40:20 - Oracle arrancado
\end{lstlisting}

\begin{lstlisting}[language=bash,label= ,caption={Ejemplo de \texttt{/home/alumno/logs/oracle.log} cuando \textbf{Oracle} se para},captionpos=b,numbers=none]
      2017-02-10-12:41:00 - Solicitud de parar Oracle
      2017-02-10-12:41:20 - Oracle parado
\end{lstlisting}

\begin{itemize}
\item Cuando no se arranque \textbf{Oracle}, porque se indica \texttt{N} en \texttt{oratab}, \textbf{systemd} intentará limpiar el servicio invocando la parada, por lo que puede que se invoque \texttt{systemctl stop} automáticamente y aparezca en el \emph{log}
\end{itemize}

\begin{lstlisting}[language=bash,label= ,caption={Ejemplo de \texttt{/home/alumno/logs/oracle.log} cuando \textbf{Oracle} no se arranca},captionpos=b,numbers=none]
      2017-02-10-12:40:00 - Solicitud de arrancar Oracle
      2017-02-10-12:40:01 - Oracle no se arranca porque /etc/oratab indica N
      2017-02-10-12:40:02 - Solicitud de parar Oracle
      2017-02-10-12:40:03 - Oracle parado
\end{lstlisting}
\end{itemize}


\begin{Aviso}
Los \emph{scripts} no cambian el fichero \texttt{/etc/oratab}, solo lo consultan.
\end{Aviso}

\begin{Aviso}
Es posible que \href{https://es.wikipedia.org/wiki/SELinux}{SELinux} no deje arrancar los \emph{scripts}. Para deshabilitarlo temporalmente, se puede usar \texttt{setenforce 0}. Para deshabilitarlo permanentemente, modificar \texttt{/etc/selinux/config}
\end{Aviso}

\section{Crea usuarios de base de datos (2 puntos)}
\label{sec:org0000006}
Crea un script de nombre \texttt{/home/alumno/scripts/nuevo-usuario-oracle.sh} que cree un nuevo usuario de oracle. Si se invoca con un número de parámetros distinto de 2, mostrará el texto de ayuda del listado \ref{lst:nuevo-usuario-ayuda}

\begin{lstlisting}[language=bash,label=lst:nuevo-usuario-ayuda,caption={Ayuda del script \texttt{nuevo-usuario-oracle.sh}},captionpos=b,numbers=none]

  Crea un usuario nuevo de oracle, con permisos connect y resource.
  Si el usuario ya existe, lo desbloquea y le cambia la contraseña.
  
  Uso: nuevo-usuario-oracle.sh <usuario> <contraseña>
\end{lstlisting}


\begin{Aviso}
En la salida del \emph{script} debe quedar claro si el usuario se crea (porque no existe), o solo es desbloqueado.
\end{Aviso}







\section{Almacena información periódicamente en la base de datos (3 puntos)}
\label{sec:org0000009}
Programa un \emph{script} para que cada minuto almacene en la tabla \texttt{DF} la información del comando \texttt{df -k}. Esta tabla (listado \ref{lst:tabla.sql}) tendrá como columnas:

\begin{itemize}
\item \texttt{hora}: Hora de lanzamiento del comando
\item \texttt{sistema}: Nombre del tipo de sistema de ficheros
\item \texttt{tamano}: Tamaño en KB del sistema de ficheros
\item \texttt{usado}: Tamaño usado, en KB
\item \texttt{montado}: Punto de montaje
\end{itemize}


\begin{lstlisting}[language=SQL,label=lst:tabla.sql,caption={Creación de la tabla \texttt{DF}},captionpos=b,numbers=none]
  create table DF(
    hora varchar(40),
    sistema varchar(40),
    tamano varchar(40),
    usado varchar(40),
    montado varchar(40)
  );
\end{lstlisting}

\begin{minipage}{\textwidth}
\begin{lstlisting}[language=bash,label=lst:df-k,caption={Ejemplo de salida del comando \texttt{df -k}},captionpos=b,numbers=none]
Filesystem     1K-blocks      Used Available Use% Mounted on
udev             4002180         0   4002180   0% /dev
tmpfs             804488     19756    784732   3% /run
/dev/sda1      237874840 183034916  42733532  82% /
tmpfs            4022440    437328   3585112  11% /dev/shm
tmpfs               5120         4      5116   1% /run/lock
tmpfs            4022440         0   4022440   0% /sys/fs/cgroup
/dev/sdb5      689521880 595546232  58926896  91% /home/windows
cgmfs                100         0       100   0% /run/cgmanager/fs
tmpfs             804488        88    804400   1% /run/user/1000
\end{lstlisting}
\end{minipage}

\begin{Aviso}
Pistas opcionales para realizar el \emph{script}:
\begin{itemize}
\item Los \emph{heredocs} pueden \href{http://superuser.com/questions/456615/how-to-pass-variables-to-a-heredoc-in-bash}{contener variables}
\item \href{https://www.cyberciti.biz/tips/processing-the-delimited-files-using-cut-and-awk.html}{Cortar columnas} con \texttt{awk}
\item Leer líneas \href{http://stackoverflow.com/questions/10929453/read-a-file-line-by-line-assigning-the-value-to-a-variable}{una por una} y meterlas en una variable:
\item \href{https://linuxconfig.org/commands-on-how-to-delete-a-first-line-from-a-text-file-using-bash-shell}{Quitar la primera línea} de la salida de \texttt{df -k} con el comando \texttt{tail}
\item El \emph{script} debería seguir los siguientes pasos:
\begin{enumerate}
\item Quitar la primera línea de la salida de \texttt{df -k}
\item Leer cada línea con \texttt{while} y \texttt{read}
\item Sacar los campos de cada línea con \texttt{awk}
\item Ejecutar una sentencia \texttt{SQL} con los datos extraidos
\end{enumerate}
\end{itemize}
\end{Aviso}

\section{Envía un correo periódicamente (2 puntos)}
\label{sec:org000000c}


\begin{itemize}
\item Programa un \emph{script} para que cada minuto envíe un correo con la información promedio del comando  \texttt{df -k}. Puedes usar como base para la consulta el listado \ref{lst:promedio.sql}
\item El correo se enviará a \texttt{alvaro@alvarogonzalez.no-ip.biz}
\item Indica tu nombre en el asunto del correo
\item El fichero tendrá \href{https://stackoverflow.com/questions/643137/how-do-i-spool-to-a-csv-formatted-file-using-sqlplus\#654723}{formato CSV}. Se debe poder abrir directamente con \textbf{excel} y visualizar su resultado en filas y columnas.
\end{itemize}



\begin{minipage}{\textwidth}
\begin{lstlisting}[language=SQL,label=lst:promedio.sql,caption={Consulta tipo para extraer información promedio},captionpos=b,firstnumber=1,numbers=left]
  select 
    sistema, avg(tamano), avg(usado), montado
  from 
    DF
  group by
    sistema, montado;
\end{lstlisting}
\end{minipage}

Para ser un servidor de correo, se necesita \href{https://docs.oracle.com/en/learn/oracle-linux-postfix/\#configure-postfix}{instalar Postfix}
\begin{lstlisting}[language=bash,label= ,caption= ,captionpos=b,firstnumber=1,numbers=left]
sudo dnf install -y postfix mailx
sudo systemctl start postfix
\end{lstlisting}


\section{Instrucciones de entrega}
\label{sec:org000000f}

\begin{itemize}
\item El ejercicio se realizará y entregará de manera individual.
\item El profesor comprobará el funcionamiento del sistema el día indicado.
\item Sube en la tarea del aula virtual un ZIP con todos los ficheros que has creado o modificado:
\begin{itemize}
\item \emph{Scripts}
\item \emph{units} de \texttt{systemd}
\item Ficheros de \texttt{cron} / \texttt{anacron}
\end{itemize}
\end{itemize}
\end{document}
